% Gauge Fixing: Optional Per-Field Toolkit
% Source: docs/gauge_fixing.md (migrated to LaTeX)
% Uses macros from preamble.tex

\section{Gauge Fixing: Optional Per-Field Toolkit}
\label{sec:gauge-fixing}

Gauge fixing simplifies equation structure for theories with gauge symmetry
(massless vectors, linearized gravity). It is \textbf{never required} --- TIDAL's
existing pipeline handles gauge-invariant theories correctly, and all
measurement quantities (energy, conversion, mixing) are gauge-invariant.

\subsection{Quick Start}
\label{sec:gauge-quick-start}

Add a \texttt{[[gauge]]} section to your \texttt{theory.toml}:

\begin{lstlisting}[style=toml]
[[fields]]
name = "A"
type = "vector"

[[gauge]]
field = "A"
type = "lorenz"
xi = 1.0              # optional, default 1.0 (Feynman gauge)
\end{lstlisting}

This adds the Lorenz gauge-fixing term $-(1/2\xi)(\partial_\mu A^\mu)^2$ to the
Lagrangian before Euler--Lagrange derivation. For Maxwell equations, this
reduces coupled equations to uncoupled wave equations for each component.

No \texttt{[[gauge]]} section means no gauge fixing (the default). Existing
examples continue to work unchanged.

\subsection{Named Presets}
\label{sec:gauge-presets}

\subsubsection{Type A --- Lagrangian Term (Added to $\lag$ Before EL Derivation)}
\label{sec:gauge-type-a}

\begin{table}[htbp]
\centering
\caption{Type A gauge-fixing presets (Lagrangian term).}
\label{tab:gauge-type-a}
\begin{tabular}{@{}llll@{}}
\toprule
Preset & Fields & Expression & Effect \\
\midrule
\texttt{lorenz} & vector & $-(1/2\xi)(\partial_\mu A^\mu)^2$ & Maxwell $\to$ uncoupled wave eqs. \\
\texttt{de\_donder} & sym.\ rank-2 & $-(1/2\xi)(\partial_a h^a{}_b - \tfrac{1}{2}\partial_b h)^2$ & Lin.\ Einstein $\to$ uncoupled waves \\
\bottomrule
\end{tabular}
\end{table}

\subsubsection{Type B --- Constraint (Applied After EOM Derivation)}
\label{sec:gauge-type-b}

\begin{table}[htbp]
\centering
\caption{Type B gauge-fixing presets (constraint).}
\label{tab:gauge-type-b}
\begin{tabular}{@{}llll@{}}
\toprule
Preset & Fields & Constraint & Effect \\
\midrule
\texttt{temporal} & vector & $A_0 = 0$ & Zeroes temporal component in all eqs. \\
\texttt{coulomb} & vector & $\nabla \cdot \mathbf{A}_\text{spatial} = 0$ & Adds transversality constraint eq. \\
\texttt{axial} & vector & $A_n = 0$ & Zeroes last spatial component in all eqs. \\
\texttt{tt} & sym.\ rank-2 & $h_{0,\mu}=0$, $h^i{}_i=0$, $\partial^i h_{ij}=0$ & TT gauge: temporal + traceless + transverse \\
\bottomrule
\end{tabular}
\end{table}

Type A presets use Wolfram builder functions in \texttt{GaugeFix.wl}.
Type B presets operate at the component level after \texttt{DecomposeToComponents},
modifying or augmenting the equation system directly.

\subsection{Custom Gauge Terms}
\label{sec:gauge-custom}

For full flexibility, use \texttt{type = "custom"} with an arbitrary Wolfram
expression and \texttt{mechanism = "lagrangian\_term"}:

\begin{lstlisting}[style=toml]
[[gauge]]
field = "A"
type = "custom"
mechanism = "lagrangian_term"
expression = "-(1/(2*xi)) * eta[a,b] CD[-a][A[-b]] eta[c,d] CD[-c][A[-d]]"
\end{lstlisting}

The \texttt{expression} field uses the same syntax as \texttt{[lagrangian].expression} and
\texttt{[[derived\_fields]].definition} --- field names, \texttt{eta}, \texttt{CD}, and chart
placeholders are all substituted automatically.

\paragraph{Note.}
Custom constraint gauges (\texttt{mechanism = "constraint"}) are not yet
supported. Use a named preset (\texttt{temporal}, \texttt{coulomb}, \texttt{axial}) for
constraint-based gauge fixing.

\subsubsection{Worked Example: Lorenz Gauge Two Ways}
\label{sec:gauge-lorenz-example}

\textbf{Using the preset:}

\begin{lstlisting}[style=toml]
[[gauge]]
field = "A"
type = "lorenz"
xi = 1.0
\end{lstlisting}

\textbf{Equivalent custom expression:}

\begin{lstlisting}[style=toml]
[[gauge]]
field = "A"
type = "custom"
mechanism = "lagrangian_term"
expression = "-(1/2) * eta[a,b] CD[-a][A[-b]] eta[c,d] CD[-c][A[-d]]"
\end{lstlisting}

Both produce the same gauge-fixed Lagrangian. The preset is more readable;
the custom form shows the underlying mechanism and serves as a template for
non-standard gauge choices.

\subsection{Two Mechanisms}
\label{sec:gauge-mechanisms}

\textbf{Type A --- Lagrangian term} (\texttt{mechanism = "lagrangian\_term"}):

\begin{itemize}
  \item An expression is added to $\lag$ \emph{before} Euler--Lagrange derivation
  \item Changes the structure of the equations of motion
  \item The gauge-fixed Lagrangian is canonicalized (\texttt{ToCanonical} + \texttt{ContractMetric})
    before EL derivation
  \item Example: Lorenz gauge adds $-(1/2\xi)(\nabla \cdot A)^2$, turning Maxwell into wave equations
  \item Presets: \texttt{lorenz}, \texttt{de\_donder}
\end{itemize}

\textbf{Type B --- Constraint} (\texttt{mechanism = "constraint"}):

\begin{itemize}
  \item Applied \emph{after} EOM derivation and component decomposition
  \item For \texttt{temporal} and \texttt{axial}: substitutes the constrained component (and all its
    derivatives) with zero via \texttt{ReplaceAll} on the component equations
  \item For \texttt{coulomb}: appends a spatial divergence constraint equation
    ($\nabla \cdot \mathbf{A}_\text{spatial} = 0$) to the equation system as a \texttt{time\_order=0} equation
  \item The original EOM structure is preserved; constraints modify or augment it
  \item For \texttt{tt} (symmetric rank-2 tensors): applies three conditions ---
    (1)~temporal: zero all $h_{0,\mu}$ components,
    (2)~traceless: substitute last spatial diagonal $h_{d-1,d-1} = -(\text{sum of other spatial diagonals})$,
    (3)~transverse: append $\partial^i h_{ij} = 0$ constraint equations for each spatial direction $j$
  \item Presets: \texttt{temporal}, \texttt{coulomb}, \texttt{axial}, \texttt{tt}
\end{itemize}

For named presets, the mechanism is inferred automatically (e.g., \texttt{lorenz}
is always \texttt{lagrangian\_term}, \texttt{temporal} is always \texttt{constraint}).

\subsection{Multi-Field Example}
\label{sec:gauge-multi-field}

Different fields can have different gauge choices, or no gauge at all:

\begin{lstlisting}[style=toml]
[[fields]]
name = "A"
type = "vector"

[[fields]]
name = "B"
type = "vector"

# Lorenz gauge on A only
[[gauge]]
field = "A"
type = "lorenz"

# B has no gauge fixing (default)
\end{lstlisting}

\subsection{Adding New Presets (Developer Guide)}
\label{sec:gauge-developer}

The process depends on the gauge mechanism.

\subsubsection{Type A (Lagrangian Term) --- 3 Steps}
\label{sec:gauge-dev-type-a}

\paragraph{Step 1:} Write a builder function in \texttt{tidal/wolfram/GaugeFix.wl}:

\begin{lstlisting}[style=mathematica]
BuildMyGaugeTerm[field_, metric_, covd_, xi_:1] := Module[
  {gaugeTerm},

  If[!NumericQ[xi] || xi <= 0,
    Throw["BuildMyGaugeTerm: xi must be positive, got " <> ToString[xi]]
  ];

  (* Your gauge-fixing expression here *)
  gaugeTerm = (* ... *);

  gaugeTerm
];
\end{lstlisting}

The function must return an \emph{expression} (a Wolfram symbolic expression),
not perform an action. The pipeline calls \texttt{AddGaugeFixingTerm[L, yourTerm]}
to inject it into the Lagrangian.

\paragraph{Step 2:} Add a registry entry in \texttt{tidal/cli/\_derive.py}:

\begin{lstlisting}[style=python]
_GAUGE_PRESETS["my_gauge"] = {
    "mechanism": "lagrangian_term",
    "builder": "BuildMyGaugeTerm",    # Wolfram function name
    "requires": "vector",             # or "tensor"
}
\end{lstlisting}

\paragraph{Step 3:} (Optional) Add validation rules in \texttt{\_validate\_gauge()}.

\subsubsection{Type B (Constraint) --- 2 Steps}
\label{sec:gauge-dev-type-b}

Type B presets operate at the component level in Python --- no Wolfram builder
function needed.

\paragraph{Step 1:} Add a registry entry (no \texttt{builder} key):

\begin{lstlisting}[style=python]
_GAUGE_PRESETS["my_constraint"] = {
    "mechanism": "constraint",
    "requires": "vector",
}
\end{lstlisting}

\paragraph{Step 2:} Add a handler in \texttt{\_wls\_gauge\_fixing\_type\_b()} in \texttt{\_derive.py}.
Use \texttt{\_type\_b\_zero\_component()} for substitution-type constraints or
\texttt{\_type\_b\_coulomb\_constraint()} as a template for differential constraints.

\subsection{FAQ}
\label{sec:gauge-faq}

\paragraph{Do I need gauge fixing?}
No. All existing TIDAL examples work without it. Gauge fixing is a
computational convenience --- it simplifies equation structure but doesn't
change the physics. Massless vectors already get implicit constraint
structure ($A_0$ becomes \texttt{time\_order=0}). Massive vectors have no gauge
freedom at all.

\paragraph{Are measurements affected by gauge choice?}
No. Energy density, conversion probability, mixing length, and spectral
quantities are all gauge-invariant physical observables.

\paragraph{Can I mix presets and custom gauges?}
Yes. Each \texttt{[[gauge]]} entry is independent. You can use \texttt{type = "lorenz"}
on one field and \texttt{type = "custom"} on another.

\paragraph{Can I mix Type A and Type B gauges?}
Yes. For example, \texttt{lorenz} on field A (Type~A) and \texttt{temporal} on field B
(Type~B) is valid. Type~A modifies the Lagrangian before EL; Type~B
modifies the component equations after decomposition.

\paragraph{Can I use gauge fixing with linearized theories?}
Yes. When \texttt{[lagrangian]} and \texttt{[linearization]} are both present, the pipeline
applies gauge fixing in the Lagrangian-first linearization path:

\begin{itemize}
  \item Type~A terms are added to $\lag^{(2)}$ after xPert expansion
  \item Type~B constraints are applied after component decomposition
\end{itemize}

This requires \texttt{[lagrangian]} --- legacy \texttt{[linearization]} without \texttt{[lagrangian]}
does not support \texttt{[[gauge]]}.

\paragraph{What happens if I gauge-fix a field that has no gauge symmetry?}
For massive fields (Proca), the mass term already breaks gauge symmetry.
Adding a gauge-fixing term is mathematically valid but physically
unnecessary and may produce unexpected equations. Validation prevents
gauge-fixing scalar fields (which never have gauge symmetry).
